\section{Resolution-dependent graph width}\label{sec:v25-graph-width}
For $0<\varepsilon\le1$ put
\begin{equation}\label{eq:v25-edge-bits}
 b_e(\varepsilon,a)=\max_{0\le\ell\le p_e}
       \left\{\frac{\ell}{2}\log_2\frac1\varepsilon
                              +\log_2 V_{e,\ell}\right\}.
\end{equation}
The zero-index term is zero and a zero positive-index volume contributes
$-\infty$. These nonnegative real profiles are not separately rounded
component memories. Define
\begin{equation}\label{eq:v25-weighted-width}
 W_G(\varepsilon,a)=\min_\pi\max_{1\le j<v}
                       \sum_{e\in\delta(S_j^\pi)}b_e(\varepsilon,a).
\end{equation}
For constant edge weights this is weighted cutwidth. The usual unweighted
parameter counts edges crossing prefixes of a vertex ordering; see
\cite{AmarilliGroz2025}. Here the weights are derived from an experiment
and change with its required predictive resolution.

\begin{corollary}[Persistent bits and graph width]\label{cor:v25-graph-bits}
Let $B_G^{*,\mathrm{av}}(\varepsilon,a)$ and
$B_G^{*,\mathrm{max}}(\varepsilon,a)$ be the smallest nonnegative
integer numbers of persistent bits achieving the respective risk at
most $\varepsilon$ in the graph experiment. Then
\begin{equation}\label{eq:v25-graph-bit-law}
 B_G^{*,\mathrm{av}}(\varepsilon,a)
       =W_G(\varepsilon,a)+O(1),\qquad
 B_G^{*,\mathrm{max}}(\varepsilon,a)
       =W_G(\varepsilon,a)+O(1).
\end{equation}
The remainders are uniform in $\varepsilon$ and calibration for the
fixed experiment. The same formula is attained by a deterministic
order chosen for that accuracy and calibration. Adaptive selection of
vertices cannot save an unbounded number of bits relative to it.
\end{corollary}
\begin{proof}
For $M=2^B$, the inequalities $\Theta_F(M,a)\le\varepsilon$ are
exactly
\[
 B\ge\max_{0\le k\le\sum_{e\in F}p_e}
       \left\{\frac{k}{2}\log_2\frac1\varepsilon+\log_2 A_{F,k}\right\}
   =\sum_{e\in F}b_e(\varepsilon,a).
\]
The last equality maximizes each independent capped index separately;
each has its feasible zero choice, so it also holds at exact collisions.
Existence of an order satisfying all these inequalities is equivalent
to $B\ge W_G(\varepsilon,a)$. The two constants in
Theorem~\ref{thm:v25-adaptive} replace $\varepsilon$ by fixed multiples.
Every affine branch in \eqref{eq:v25-edge-bits} has slope between zero
and $p_e/2$ in $\log_2(1/\varepsilon)$. Such replacements therefore
change every cut sum, its maximum and its minimum by a bounded amount.
The same zero-term definition can be used when the replaced argument
exceeds one. Integer rounding costs at most one more bit; unused labels
may be padded. Apply the upper and lower comparisons to obtain both
formulas. Proposition~\ref{prop:v25-fixed-order} supplies the matching
deterministic construction.
\end{proof}

At a fixed calibration define the integer edge weights
\[
 d_e(a)=\min\{n_e(r_e-1),|m_eA_e|-1\},\qquad
 \Gamma_G(a)=\min_\pi\max_j\sum_{e\in\delta(S_j^\pi)}d_e(a).
\]
\begin{corollary}[The fixed-calibration exponent]\label{cor:v25-cutwidth}
If $E\ne\varnothing$, both optimal graph risks at a fixed calibration
have order $M^{-2/\Gamma_G(a)}$. For the binary experiments
$A_e=\{0,1\}$ and $n_e=m_e=1$, this is $M^{-2/\operatorname{cw}(G)}$,
where $\operatorname{cw}(G)$ is ordinary cutwidth.
\end{corollary}
\begin{proof}
At fixed calibration the last positive volume index of edge $e$ is
$d_e(a)$, and every preceding index is positive. At cut $F$ the last
positive index of $A_{F,k}$ is consequently $\sum_{e\in F}d_e(a)$.
The finite positive coefficients in its profile give order
$M^{-2/\sum_{e\in F}d_e(a)}$ for $M\ge1$ and nonempty $F$.
Taking the maximum over cuts, the minimum over finitely many orders,
and applying Theorem~\ref{thm:v25-adaptive} proves the assertion.
Here coefficient comparisons may depend on that fixed calibration.
Binary two-trial edges have $d_e=1$, which gives ordinary cutwidth.
\end{proof}

\subsection{A change of the optimal order with accuracy}
Unlike unconstrained serial schedules, a graph need not admit one
order with bounded additive excess uniformly over the collision
calibrations and resolutions considered below. At each fixed
calibration, Theorem~\ref{thm:v30-menu-finite} supplies a bounded-excess
order on the entire accuracy interval. The next example uses only
three independent latent coordinates.

\begin{corollary}[A multiscale ordering transition]\label{cor:v25-star}
There is a fixed three-edge star experiment, with one known calibration
$0\le\theta\le\theta_0$, using exactly $24$ raw trials, for which
both optimal bit criteria satisfy, as $\theta\downarrow0$,
\begin{equation}\label{eq:v25-star-bits}
 B_G^*(\theta,\theta)=\tfrac72\log_2(1/\theta)+O(1),\qquad
 B_G^*(\theta^4,\theta)=16\log_2(1/\theta)+O(1).
\end{equation}
For each sufficiently small $\theta$, no single deterministic vertex
order is within bounded additive bits of the optimum at both indicated
accuracies, uniformly in $\theta$. More precisely, every order has an
excess of at least $\tfrac12\log_2(1/\theta)-C$ at one of the two
accuracies, for a fixed $C$. The order may depend on $\theta$ in this
assertion, but must be the same at the two accuracies.
\end{corollary}
\begin{proof}
Label the star's three edges $1,2,3$. For edges $1$ and $3$ take
$A_e=\{0,1\}$, the cells
$k_0=1/2+t/4$, $k_1=1/2-t/4$, and respectively
$(n_1,m_1)=(7,7)$ and $(n_3,m_3)=(1,1)$.
For edge $2$ take $A_\theta=\{0,1,2+\theta\}$,
$(n_2,m_2)=(5,3)$, and
\[
 k_0=\frac13+\frac{t+t^{2+\theta}}{12},\qquad
 k_1=\frac13-\frac{t}{12},\qquad
 k_2=\frac13-\frac{t^{2+\theta}}{12}.
\]
The cells sum to one, have fixed coefficient matrices of full column
rank, and are bounded below by $1/4$. Use any fixed full-support prior
in each edge and fixed admitted command cubes. Take, for instance,
$\theta_0=1/16$. The total raw trial count is $14+8+2=24$.

Write $L=\log_2(1/\varepsilon)$ and $U=\log_2(1/\theta)$.
The first and third edges give
\[
 b_1=\tfrac72L+O(1),\qquad b_3=\tfrac12L.
\]
For the second edge the formal cap is $p_2=9$. Its nine positive
three-fold sums, before the fixed normalization, are
\[
 1;\quad 2,2+\theta;\quad 3,3+\theta;\quad
 4+\theta,4+2\theta;\quad 5+2\theta;\quad 6+3\theta.
\]
There are six uniformly separated groups, three of them pairs with
gap $\theta$. A subset of size $\ell>6$ must contain at least
$\ell-6$ complete pairs, and this minimum is achieved by using all
six groups first. For $\ell\le6$, choose at most one point per group.
All other pair distances are bounded above and below by positive
constants. Thus
\[
 V_{2,\ell}\asymp\theta^{(\ell-6)_+}\quad(1\le\ell\le9),\qquad
 b_2=\max\{3L,\tfrac92L-3U\}+O(1).
\]
The constants are independent of $L$ and $\theta>0$. At $\theta=0$
the last three volumes vanish and the original graph theorem still
applies; logarithms here are used only for positive $\theta$.

For arbitrary nonnegative star weights $w_1,w_2,w_3$, an order places
a subset of leaves before the centre and its complement after it.
The maximum cut weight is
$\max\{\sum_{i\in A}w_i,\sum_{i\notin A}w_i\}$: the successive
weights increase to the first sum and decrease from the second sum.
Hence minimization is a two-part partition problem. With three positive
weights an optimal partition isolates a largest weight. To see this,
write $w_1\ge w_2\ge w_3$. Isolating $w_1$ costs
$\max(w_1,w_2+w_3)\le w_1+w_3$, the cost of isolating $w_2$;
isolating $w_3$ or taking an empty part is no better.

At $\varepsilon=\theta$, the weights are
$(7U/2,3U,U/2)+O(1)$. Isolating edge $1$ costs $7U/2+O(1)$,
whereas isolating edge $2$ costs $4U+O(1)$, isolating edge $3$ costs
$13U/2+O(1)$, and an empty part costs $7U+O(1)$.
At $\varepsilon=\theta^4$ the weights are
$(14U,15U,2U)+O(1)$. Isolating edge $2$ costs $16U+O(1)$,
isolating edge $1$ costs $17U+O(1)$, and the other choices cost at
least $29U+O(1)$. Corollary~\ref{cor:v25-graph-bits} proves
\eqref{eq:v25-star-bits}. Its fixed-order version follows by the same
inversion from Proposition~\ref{prop:v25-fixed-order}. Comparing the
four possible cut partitions proves the stated uniform excess for
every fixed order. The bounded remainders can be made simultaneous
because the number of orders is finite.
\end{proof}

The ordering transition is not determined by the generic rank of the
second edge. Its three unresolved collision scales determine when
that edge outweighs the first. The adaptive converse shows that the
best resolution-dependent graph layout already gives the optimal
order of memory for the stated experiment, even against measurable
history-dependent choices. Independence of the edge priors, exogenous
commands and the specified batch boundaries are substantive hypotheses;
no conclusion for arbitrary correlated latent coordinates or different
checkpoint protocols is implicit.
